\section{Introduction}

The ability to accurately predict the future is a foundational element of human intelligence and a prerequisite for effective planning in modern society. From managing energy consumption in smart grids to navigating complex geopolitical shifts, the stakes of forecasting are high. Traditionally, this field has been split into two distinct regimes: numerical time-series forecasting, which relies on high-volume data and statistical patterns, and event-based forecasting, which requires abstract reasoning over sparse, qualitative information.

{\bf Why does this matter?} In the numerical regime, human cognition is poorly suited for processing the vast arrays of data generated by modern sensors and markets. Conversely, in the event-based regime, human experts---often termed ``superforecasters''---have historically set the gold standard by leveraging heuristic-based reasoning. However, as the world becomes increasingly data-rich and complex, the need for a unified forecasting framework that can excel across both high-data and low-data regimes has become urgent.

{\bf What is the gap?} Despite the rapid advancement of Large Language Models (LLMs), existing research typically isolates these two forecasting domains. While work like \timellm \cite{jin2024timellm} and \onefitsall \cite{zhou2023onefitsall} has shown that LLMs can be reprogrammed for numerical tasks, and benchmarks like \forecastbench \cite{karger2025forecastbench} have evaluated event-based reasoning, a unified assessment that directly contrasts LLM performance against human benchmarks across both extremes is missing.

In this work, we hypothesize that LLMs can match or surpass human performance in both short-term high-data and long-term low-data scenarios. We evaluate \methodname on the \ettm dataset for numerical time-series forecasting and on a subset of \forecastbench for event-based probabilistic forecasting. Our approach avoids complex domain-specific fine-tuning, focusing instead on the zero-shot capabilities of foundational models.

{\bf Preview of results.} Our experiments reveal that \methodname achieves a Brier score of 0.246 on long-term event forecasting, outperforming human superforecasters (0.287) by 14.3\% and remaining comparable to the aggregated public crowd (0.243). In the numerical regime, although the LLM slightly underperforms simple statistical baselines like moving averages (MSE 3.321 vs 3.988), it demonstrates a zero-shot capacity to process high-volume numerical data that is fundamentally impossible for human manual forecasting.

We make the following contributions:
\begin{itemize}[leftmargin=*,itemsep=0pt,topsep=0pt]
    \item We conduct a unified evaluation of LLMs on two forecasting extremes: short-term high-data numerical and long-term low-data event-based.
    \item We empirically demonstrate that \methodname surpasses human expert superforecasters in probabilistic event prediction.
    \item We provide a performance analysis of zero-shot LLMs on high-volume numerical data, establishing a baseline for general-purpose models in data-intensive regimes.
\end{itemize}

The remainder of this paper is organized as follows: Section 2 reviews related work in LLM-based forecasting; Section 3 describes our methodology; Section 4 presents our experimental results; Section 5 discusses limitations and implications; and Section 6 concludes.
